\documentclass[10pt]{article}
\usepackage[margin=1in]{geometry}
\usepackage{booktabs}
\usepackage{array}
\usepackage{amsmath,amssymb}
\usepackage{graphicx}
\usepackage[hidelinks]{hyperref}
\usepackage{enumitem}
\usepackage{listings}
\usepackage{xcolor}
\lstset{basicstyle=\ttfamily\small,breaklines=true,frame=single,
  framerule=0.2pt,columns=fullflexible,keepspaces=true}
\setlist{itemsep=1pt,topsep=3pt}
\emergencystretch=2.5em
\hbadness=10000

\title{Supplementary Material for\\
\emph{Look Before You Look: Auditing Perception Utility and Evaluation
Sensitivity in Multimodal Agents}}
\author{Anonymous Author(s)\\AAMAS 2027 Submission}
\date{}

\begin{document}
\maketitle

\tableofcontents

\section*{Overview}

This document provides implementation details deferred from the main paper
(exact crop coordinates, prompts, filtering rules, fallbacks, OCR, and scoring
details; Sections~\ref{sec:budget}--\ref{sec:features}), frozen policy and
controller specifications (Section~\ref{sec:policies}), the preregistration
index and deviations log (Section~\ref{sec:prereg}), and additional analyses
supporting the paper's three-level distinction among unrestricted structural
opportunity, action-matched counterfactual opportunity, and frozen prospective
policy performance (Section~\ref{sec:results}). Section~\ref{sec:results}
also reports the blinded human-adjudication analysis and the nonlinear and
representation--head prospective follow-ups. Sections~\ref{sec:infra}--\ref{sec:zip}
document infrastructure, data separation, and the accompanying
reproducibility artifacts.

All automatically scored quantities are recomputable offline from the
released result files because raw generations are preserved for every
rollout. Human adjudication is provided separately through the frozen
annotation sample, blinded labels, and adjudication summaries.

% ======================================================================
\section{Visual-Token Budget Accounting}
\label{sec:budget}

Visual cost is measured in Qwen2.5-VL visual tokens, where one token covers
a $28\times28$ pixel block after smart resizing (14-pixel patches with
$2\times2$ merging):
\[
\text{tokens}(h,w) = \left\lfloor \bar h/28 \right\rfloor \cdot
\left\lfloor \bar w/28 \right\rfloor,
\qquad (\bar h,\bar w) = \mathrm{smart\_resize}(h,w),
\]
where $\mathrm{smart\_resize}$ mirrors the official
\texttt{qwen\_vl\_utils.smart\_resize}: aspect ratio is preserved, both
sides are snapped to multiples of 28, and the total area is clamped to
$[\text{min\_pixels},\,\text{max\_pixels}]$ with
$\text{max\_pixels} = B \cdot 28^2$ for a budget of $B$ tokens
(code: \texttt{ppu/imaging.py}).

We resize images ourselves and transmit the resized bytes (JPEG, quality
92, no chroma subsampling) rather than passing server-side pixel limits, so
the realized budget is identical across backends and replays and is
auditable from the logs. Every rendered view is logged together with its
realized token count and a SHA-1 hash of its pixels, which lets the audit
verify that distinct crop candidates are distinct observations.

\paragraph{Per-dataset budgets.}
Coarse / \textsc{HighRes} visual-token targets are 256/1{,}280 for
MME-RealWorld-Lite, 192/512 for TextVQA, and 64/192 for VQAv2. Crop views
are re-encoded at 256 tokens on all datasets. The full-information
reference view (used only for diagnostics, never as a policy action) is
capped at 2{,}560 tokens. A sample enters the audit only if its native view
costs at least $4\times$ the coarse budget; otherwise \textsc{Crop} and
\textsc{HighRes} cannot add information (inclusion rule).

Views below $4$ tokens are rejected as unreadable
($\text{min\_pixels}=4\cdot28^2$), and crop boxes that map to fewer than
8 pixels on a side raise an error rather than silently degrading.

% ======================================================================
\section{Perception Actions: Exact Specifications}
\label{sec:actions}

\subsection{Crop candidate construction ($K=8$, constant cardinality)}

Every evaluated state has exactly eight crop candidates, but the
construction depends on the experimental phase. Constant cardinality is
important because the unrestricted oracle maximizes over crop candidates;
varying the number of candidates would change the amount of
selection-on-noise opportunity.

\paragraph{Initial structural audit.}
The initial structural audit (Gate-0, both readers) uses five deterministic
multiscale grid boxes together with up to three frozen question-conditioned
proposals. Proposal generation, parsing, filtering, and deterministic
padding follow the procedure described below.

\paragraph{Held-out and prospective collections.}
All held-out and prospective collections remove VLM-generated proposals and
use the same eight fixed boxes for every image: the five grid boxes below
plus three smaller deterministic fallback boxes. Thus the deployed selectors,
the state-independent mixtures, and the action-matched oracle operate over
an identical fixed crop family.

\paragraph{Deterministic multiscale grid (5 boxes).}
Normalized $(x_0, y_0, x_1, y_1)$ coordinates:
\begin{center}
\begin{tabular}{ll}
\toprule
Region & Box \\
\midrule
Center       & $(0.25,\,0.25,\,0.75,\,0.75)$ \\
Top-left     & $(0.00,\,0.00,\,0.55,\,0.55)$ \\
Top-right    & $(0.45,\,0.00,\,1.00,\,0.55)$ \\
Bottom-left  & $(0.00,\,0.45,\,0.55,\,1.00)$ \\
Bottom-right & $(0.45,\,0.45,\,1.00,\,1.00)$ \\
\bottomrule
\end{tabular}
\end{center}
The sequential experiment's frozen $\mathrm{CROP}_{\mathrm{TR}}$ action is
the top-right grid box above; its first-step center crop is the center grid
box.

\paragraph{Fixed fallback boxes for held-out and prospective collections.}
The three additional fixed boxes are
\[
(0.35,\,0.35,\,0.65,\,0.65),\quad
(0.05,\,0.05,\,0.35,\,0.35),\quad
(0.65,\,0.65,\,0.95,\,0.95),
\]
deliberately smaller than the grid so that they probe a different scale.
Together with the five grid boxes, these define the fixed eight-crop family
used in all held-out and prospective one-step evaluations (crop budget 256
tokens each, coarse overview retained).

\paragraph{Question-conditioned proposals in the initial structural audit only.}
A single proposal generation is requested from the frozen reader with the
prompt in Section~\ref{sec:prompts:proposal}. Its cost is charged to every
arm that considers \textsc{Crop} at all, including random and oracle arms.
Parsing is tolerant but never creative:
\begin{itemize}
\item Each output line must contain at least four numbers; lines that do
  not are dropped, not repaired.
\item Qwen2.5-VL is grounded in absolute pixels and frequently mixes pixel
  and fractional coordinates in one line. Each coordinate is therefore
  interpreted independently: a value with $|v|\le 1$ is read as a fraction,
  otherwise it is divided by the width or height of the view the model
  actually saw. Every parsed box logs its raw values, per-coordinate
  interpretation mode, and the reason for any rejection, so the decision is
  auditable per line. (An earlier all-fractional parser silently dropped
  $\sim$70\% of proposal calls; the mixed parser and this failure are
  disclosed in the code comments.)
\item Filters, applied in order: out-of-range after interpretation
  (tolerance $\pm 0.01$); degenerate or below 1\% of image area; area
  above $0.95$ (near-duplicate of the full view); IoU $> 0.9$ with an
  already accepted box. At most 3 proposals are kept.
\end{itemize}

\paragraph{Padding to $K=8$ in the initial structural audit.}
Grid and proposal boxes are merged (IoU $>0.9$ deduplication, grid first).
If fewer than 8 remain, deterministic fallback boxes are appended:
$(0.35,0.35,0.65,0.65)$, $(0.05,0.05,0.35,0.35)$,
$(0.65,0.65,0.95,0.95)$ --- deliberately smaller boxes so a fallback probes
a different scale rather than duplicating the grid. If still short (heavy
overlap), centered square boxes of side $0.45, 0.40, 0.35, \ldots$
(step $0.05$, floor $0.12$) are added. The candidate set is thus a
deterministic function of the sample and the single proposal generation.

\paragraph{Crop semantics.}
Crops are taken from the \emph{original} native-resolution image, then
re-encoded at the crop budget; cropping the coarse view would make
\textsc{Crop} an upsampling artifact rather than an information-acquiring
action. A crop view is \emph{added} to the coarse overview (global context
is retained), with a view label naming its normalized coordinates.

\paragraph{Deployable \textsc{Crop}-family semantics.}
For the deployed one-step selectors, \textsc{Crop-family} does not denote
outcome-dependent selection of the best crop. Its utility and cost are the
uniform expectation over the fixed eight-crop family. Operationally,
selecting \textsc{Crop-family} corresponds to sampling one of the eight
fixed crops uniformly at deployment. Because all eight crop outcomes are
executed in the offline audit, the reported family-level outcome and
visual-token cost are the exact expectation of this realizable randomized
policy. The unrestricted structural oracle may instead select among
individual crop outcomes and is therefore a distinct counterfactual
benchmark (Section~\ref{sec:results:deporacle}).

\subsection{Control (dummy) actions and the luck floor}

To price in ``the reader was simply asked again,'' the audit includes a
dummy arm of \emph{matched cardinality} (11 actions, same as the real
action set):
\begin{itemize}
\item \textbf{Size-matched random crops}: same count and same
  width/height distribution as the real candidates, position drawn
  uniformly from a seed derived deterministically from the sample id
  (\texttt{ppu-dummy::\{sample\_id\}}). Because area and aspect ratio are
  matched, ``question-conditioned beats random'' cannot be a crop-size
  effect in disguise.
\item \textbf{Re-encode dummies}: the coarse view re-rendered at the same
  budget with a different resampling filter (Lanczos, or bilinear plus a
  JPEG-90 round trip), i.e.\ cosmetically different pixels with no new
  information. Two distinct variants exist so the dummy pool holds 11
  \emph{distinct} draws.
\end{itemize}
The dummy oracle serves as a diagnostic control for selection-on-noise
(the ``luck floor'' of the structural audit) rather than as the comparator
defining the reported cross-replay headroom $H_{\mathrm{XR}}$, which is the
cross-replay oracle minus the cross-fitted state-independent best fixed
arm.
For question-conditioned localization specifically, size-matched random
boxes with identical area and aspect ratio (position-only randomization)
are the comparison class.

\subsection{OCR as an external tool}

\begin{itemize}
\item Engine: RapidOCR (ONNX runtime, CPU). The OCR engine is deliberately
  \emph{not} the VLM: if the reader transcribed the image itself, OCR and
  \textsc{HighRes} would be the same action and action-diversity results
  would be a prompt-wording artifact.
\item OCR reads the \emph{original} native-resolution image, not the coarse
  view --- extracting text the reader cannot resolve is the point of the
  tool.
\item Serialization: detected lines are sorted into reading order (by
  rounded $y_0$, then $x_0$), truncated to 30 lines, and rendered one line
  per row as \verb|[x0,y0,x1,y1] text| with normalized coordinates at two
  decimals. An empty result is rendered as \verb|(no text detected)|.
\item Cost accounting: engine latency and the tool call are recorded as
  cost. Campaigns use a pre-pass cache (each image OCRed once across CPU
  workers), but the \emph{recorded} latency is the engine's real cost, not
  the cache lookup's, so the cost column reports what a deployed agent
  would pay. OCR output is an observation, produced only after the action
  is chosen; nothing upstream of the action choice may read it.
\end{itemize}

\subsection{\textsc{HighRes} and \textsc{Stop}}

\textsc{HighRes} replaces the coarse overview with a higher-resolution full
view at the dataset's \textsc{HighRes} budget. \textsc{Stop} answers from
the coarse view alone. In the sequential experiment, \textsc{HighRes} as a
second step replaces the coarse overview while \emph{retaining} evidence
already acquired by the first action.

% ======================================================================
\section{Prompts (Verbatim)}
\label{sec:prompts}

One template serves every arm; the only thing that varies across actions is
the observation block. If actions had their own instructions, action
comparisons would partly measure prompt engineering. The prompt version is
logged with every rollout and is part of the generation-cache key.

\paragraph{System prompt.}
\begin{lstlisting}
You are answering questions about images. You may be shown a low-resolution overview, additional zoomed-in views, or text extracted from the image. Use only what you are shown. Answer as briefly as possible.
\end{lstlisting}

\paragraph{Open-ended instruction, original (\texttt{p0.1}).}
\begin{lstlisting}
Answer the question with a short phrase (a few words at most). End your reply with a line of the form:
Answer: <your answer>
\end{lstlisting}

\paragraph{Open-ended instruction, one-word control (\texttt{p0.2-oneword}).}
Used only in the answer-format intervention arms:
\begin{lstlisting}
Answer with a single word or number. For yes/no questions answer exactly 'yes' or 'no'. Do not write a sentence. End your reply with a line of the form:
Answer: <your answer>
\end{lstlisting}

\paragraph{Multiple-choice instruction.}
\begin{lstlisting}
Answer with the letter of the correct option. End your reply with a line of the form:
Answer: <letter>
\end{lstlisting}

\paragraph{User-turn assembly.}
The user turn concatenates, in order and separated by blank lines: (i) a
view inventory, e.g.\ ``\texttt{Views provided, in order: coarse@192;
zoomed view of region 0.45,0.00,1.00,0.55.}''; (ii) if OCR was invoked,
``\texttt{Text extracted from the image by an OCR tool:}'' followed by the
serialized block of Section~2.3; (iii) the question as served under the
current replay (options inlined as \texttt{A.}--\texttt{E.} lines for
multiple choice); (iv) the applicable instruction above. View labels
describe what an observation \emph{is} (coordinates are part of the action
specification), never the content of an observation the agent has not paid
for.

\paragraph{Crop-proposal prompt.}
\label{sec:prompts:proposal}
\begin{lstlisting}
Question: {question}

Look at the image and the question. Name up to 3 regions you would need to inspect at higher resolution to answer it. Reply with one region per line, each as four numbers between 0 and 1: x0,y0,x1,y1 (left, top, right, bottom as fractions of the image). No other text.
\end{lstlisting}

% ======================================================================
\section{Replay Construction}
\label{sec:replays}

Decoding is greedy, so re-running a request returns the same string; a
replay therefore cannot and does not measure decoding noise. What threatens
the audit is \emph{luck} --- an action that happens to push a 5-way multiple
choice onto the right letter, or a phrasing the reader latches onto. R2
changes what is asked, never how it is sampled. Both replays are
deterministic functions of the sample id and are reproducible from the
manifest alone.

\paragraph{Multiple choice (MME-RealWorld).}
R2 permutes the answer options with a seeded RNG
(\texttt{ppu-r2::\{sample\_id\}::7}) and follows the gold answer to its new
letter. If the drawn permutation is the identity, the order is reversed
instead --- a permutation that changes nothing is not a control. Scoring
under R2 uses the moved gold letter.

\paragraph{Open-ended (TextVQA, VQAv2).}
R2 applies a fixed, meaning-preserving, syntax-level rewrite --- no LLM
sampling. Rules are applied in order, first match wins
(case-insensitive, anchored at both ends):

\begin{center}
\small
\begin{tabular}{ll}
\toprule
Pattern & Rewrite \\
\midrule
\texttt{what is \ldots?} & \texttt{Looking at the image, what is \ldots?} \\
\texttt{what are \ldots?} & \texttt{Looking at the image, what are \ldots?} \\
\texttt{what \ldots?} & \texttt{Based on the image, what \ldots?} \\
\texttt{who/where/when/which/how \ldots?} & \texttt{Based on the image, \emph{wh-word} \ldots?} \\
\texttt{is/are/does/do \ldots?} & \texttt{Based on the image, is/are/does/do \ldots?} \\
(no match) & \texttt{Based on the image: \ldots} \\
\bottomrule
\end{tabular}
\end{center}

The rewrite set is deliberately small and mechanical: a paraphrase that
changed difficulty would surface as fake instability in the label-reliability
numbers. Multi-replay robustness checks (replays 3--5, prefix variants) are
summarized in Section~\ref{sec:results:multireplay}.

% ======================================================================
\section{Scoring Rules}
\label{sec:scoring}

\paragraph{Answer extraction.}
The answer span is pulled from the raw generation in a fixed, deliberately
forgiving order: (1) an \verb|<answer>...</answer>| tag; (2) a line
matching \verb|(final )?answer[:-]|; (3) the last non-empty line. A parser
that silently preferred one surface form would bias whichever action makes
the model chatty. For multiple choice, the first standalone letter A--E in
the extracted span (falling back to the full text) is used.
\texttt{parse\_ok=False} (no answer found) scores 0 but is tracked
separately: a rising parse-failure rate indicates a pipeline bug, not a
model result.

\paragraph{Strict scorers (primary).}
No LLM judge anywhere; every metric is deterministic and recomputable
offline from the stored generations.
\begin{itemize}
\item \texttt{vqa\_acc} (TextVQA, VQAv2): official VQA accuracy.
  Normalization follows the official evaluator (lowercasing, punctuation
  stripping with number-comma and decimal-period handling, contraction
  restoration, word-to-digit mapping for zero--ten, article removal); the
  score is $\min(m/3,\,1)$, where $m$ is the number of exact matches among
  the 10 references.
\item \texttt{mcq} (MME-RealWorld): extracted letter vs.\ gold letter under
  the current replay's option order (first character, case-insensitive).
  Because scoring is letter-level, lenient and strict scoring coincide on
  MME by construction.
\item Binary quantities (\textsc{Help}/\textsc{Harm}/DMG, strict labels)
  threshold soft scores at $0.5$.
\end{itemize}

\paragraph{Format-tolerant (lenient) rescoring (secondary).}
As defined in the main paper: for normalized prediction $p$ and reference
$g_j$, a reference counts as matched if $g_j = p$, or $g_j$ equals the
first word of $p$, or $g_j$ occurs in $p$ as a whitespace-bounded
substring; the score is $\min(\sum_j h(g_j,p)/3,\,1)$. Lenient evaluation
changes only the mapping from an existing generation to its score ---
generations, observations, and replays are identical bytes.

\paragraph{Confidence.}
Primary confidence is the mean selected-token log-probability over the
complete \textsc{Stop} reply, averaged across the two \textsc{Stop}
replays. Answer-span mean, span-minimum, and first-token variants are
computed from the already-logged per-token probabilities
(Section~\ref{sec:results:spanconf}).

% ======================================================================
\section{Diagnostic and Selector Features}
\label{sec:features}

The frozen v2.2 feature set (also ``F0'' in the selector experiments)
contains 16 inexpensive pre-action features
(\texttt{ppu/gate\_features.py}). By design it contains \emph{no}
standalone confidence term: an earlier version learned a single global
direction (``confident, therefore look'') that fit VQAv2 --- where
confidence is inverted under the original prompt --- and destroyed TextVQA.
Confidence therefore enters only through interactions, so the model can
decide per answer form whether the uncertainty is worth believing.
Deliberately absent: dataset id, image resolution, native token count ---
features that name the domain rather than the state.

\begin{center}
\small
\begin{tabular}{lll}
\toprule
Group & Feature & Definition \\
\midrule
Answer surface & \texttt{ans\_chars} & characters in the provisional \textsc{Stop} answer \\
& \texttt{ans\_tokens} & generated tokens in the answer \\
& \texttt{ans\_words} & whitespace words in the answer \\
& \texttt{is\_numeric} & answer is a bare number \\
& \texttt{is\_yesno} & answer starts with yes/no \\
& \texttt{has\_hedge} & hedge regex\footnotemark\ fires on the answer \\
\midrule
Question surface & \texttt{q\_chars}, \texttt{q\_words} & question length \\
& \texttt{q\_is\_howmany} & contains ``how many'' \\
& \texttt{q\_is\_what} & starts with ``what'' \\
& \texttt{q\_is\_isare} & starts with is/are/does/do/was/were \\
\midrule
Confidence & \texttt{conf\_x\_yesno} & $\mathrm{conf\_rank}\times\texttt{is\_yesno}$ \\
(interaction only) & \texttt{conf\_x\_numeric} & $\mathrm{conf\_rank}\times\texttt{is\_numeric}$ \\
& \texttt{conf\_x\_other} & $\mathrm{conf\_rank}\times\texttt{is\_other}$ \\
& \texttt{conf\_x\_len} & $\mathrm{conf\_rank}\times\min(\texttt{ans\_chars}/20,2)$ \\
& \texttt{conf\_x\_hedge} & $\mathrm{conf\_rank}\times\texttt{has\_hedge}$ \\
\bottomrule
\end{tabular}
\end{center}
\footnotetext{\texttt{maybe|possibly|probably|appears|seems|likely|unclear|
cannot|can't|unable|not sure|hard to|difficult|blurry|unknown|unreadable|
too small} (word-bounded, case-insensitive).}

$\mathrm{conf\_rank}\in[0,1]$ is the empirical CDF of the \textsc{Stop}
mean-token log-probability, fitted on training states only (0 = most
confident). Features are standardized with training-set statistics.
The binary v2.2 diagnostic uses class-balanced logistic regression on these
features; the multi-action selector uses per-family ridge value heads on
the same features (Section~\ref{sec:policies}).

% ======================================================================
\section{Frozen Policies, Baselines, and Artifacts}
\label{sec:policies}

Each learned component is trained only on its designated training pool and
frozen before outcomes from its own prospective evaluation set are observed.
Freeze artifacts are stored (pickled) with JSON manifests recording the
relevant weights, scaler statistics, confidence ECDF, feature list, and
seed. Freeze artifacts and their SHA-256 prefixes are included in the
ZIP (\texttt{models/}).

\paragraph{One-step selector v2 (primary prospective test).}
Chosen by five-fold development cross-validation from a pre-specified
10-candidate grid (2 feature sets $\times$ 5 head classes), selection
criterion TextVQA policy value at $1/2\times$ budget, tie-break
$1/4\times$. Winner: F0 features + ridge regression ($\alpha=0.3$). Three
heads predict the replay-averaged score increment of each family
(\textsc{Crop}, \textsc{OCR}, \textsc{HighRes}) relative to \textsc{Stop};
the \textsc{Crop} target and cost use the deployable uniform-grid family
expectation (no outcome-dependent region selection). Decision rule:
$\arg\max_f [\widehat\Delta_i(f) - \lambda_B c_i(f)]$ with zero predicted
gain and cost for \textsc{Stop}; $\lambda_B$ selected on development data
inside the frozen artifact's $\lambda$ grid, per budget. Artifact:
\texttt{selector\_frozen\_v2} (sha256[:16] \texttt{d4e495063ac56b59}).

\paragraph{State-independent randomized mixtures.}
For each budget, an LP over the fixed arms \{\textsc{Stop}, each
grid/fallback crop box, uniform \textsc{Crop}, \textsc{OCR},
\textsc{HighRes}\} maximizes the mean fixed-arm outcome subject to one
expected-token constraint and a simplex constraint; with one linear
constraint the optimum mixes at most two arms (concave envelope). Three
comparator strengths are reported: in-sample (weights fitted on the
evaluation sample --- a baseline-favoring upper bound), cross-fitted
(2-fold), and deployable (fitted on development data, applied unchanged).
The frozen deployable TextVQA mixture is:
$1/4\times$: $0.876\cdot\text{CROP-center} + 0.124\cdot\text{OCR}$;
$1/2\times$: $0.569\cdot\text{CROP-center} + 0.431\cdot\text{CROP-top-right}$;
$1\times$: $1.0\cdot\text{CROP-top-right}$.

\paragraph{Hidden-state selector v3.}
Six candidate configurations compared on development data before any fresh
outcome was observed; frozen winner: the reader's \textsc{Stop}-time hidden
representation with ridge value heads ($\alpha=30$), same
value-minus-cost decision rule. Artifact: \texttt{selector\_frozen\_v3}
(sha256[:16] \texttt{4fd1fa4f}). Development diagnostics recorded before
prospective collection flagged both weaker development performance than F0
and weak $\lambda$-budget controllability; both materialized prospectively
(Section~\ref{sec:results:v3}).

\paragraph{F0 nonlinear-head prospective follow-up.}
A separately preregistered follow-up (\texttt{PREREG\_STRONG\_\allowbreak{}CONTROLLER},
09-20) tests whether the prospective gap is specific to the ridge value
head. Ridge, XGBoost, and MLP heads use the same F0 decision-time
representation, utility targets, costs, and value-minus-cost decision rule.
Training uses a frozen 1,350-state pool assembled from previously analyzed
development and evaluation states (Gate-0 200, held-out 250, prospective
600, v3 300); no state from the new 400-state evaluation set is included.
Per head class one configuration is chosen from a disclosed grid by nested
five-fold cross-validated policy value at $1/2\times$; $\lambda_B$ per
class is chosen on pooled out-of-fold predictions. The preregistered primary
comparison is F0/XGBoost minus the frozen state-independent mixture at the
strict $1/2\times$ operating point; other heads, budgets, and
format-tolerant analyses are secondary. Artifact:
\texttt{strong\_controllers\_frozen} (sha256[:16] \texttt{5c60ccab3f7d08b8}).

\paragraph{Representation--head $3\times3$ prospective test.}
The final prospective controller test (\texttt{PREREG\_HIDDEN\_\allowbreak{}CONTROLLER},
09-20) crosses three decision-time representations---F0, the
3,584-dimensional \textsc{Stop}-time hidden representation H, and
H+F0---with ridge, XGBoost, and MLP value heads. All nine configurations use
the same 1,750-state training pool (the 1,350 states above plus the
400-state nonlinear-head evaluation set) and the same family-level
value-minus-cost decision rule; parameters and per-budget $\lambda_B$ are
frozen before evaluation on a new 400-state TextVQA set, and the primary
nonlinear cell (H/MLP, by the same cross-validation rule) is recorded in the
frozen manifest before collection. The strict $1/2\times$ H/MLP-minus-mixture
comparison is preregistered as primary; other representation--head cells,
budgets, and format-tolerant analyses are secondary. Artifact:
\texttt{hidden\_controllers\_frozen} (sha256[:16] \texttt{ed93f19437264e56}).

\paragraph{Sequential controller v1 (original, preregistered test).}
Two post-first-observation features: whether the normalized first-step
answer differs from the original \textsc{Stop} answer, and first-step
answer length. Two ridge heads ($\alpha=1.0$) predict the incremental
value of $\mathrm{CROP}_{\mathrm{TR}}$ and \textsc{HighRes} over stopping;
decision $\arg\max(0, \widehat\Delta_{\mathrm{TR}},
\widehat\Delta_{\mathrm{HI}})$; no cost term. Strict-trained and
lenient-trained variants both frozen
(\texttt{seq\_controller\_frozen}, sha256[:16] \texttt{cffa1043f455f7d1}).
The frozen state-independent continuation mixture, fitted on the 100-state
development audit, places all mass on $\mathrm{CROP}_{\mathrm{TR}}$
(i.e., always-$\mathrm{CROP}_{\mathrm{TR}}$).

\paragraph{Sequential controller v2 (disclosed cost-aware follow-up).}
Seven history features (including confidence before/after the first
observation and their change), ridge continuation-value heads, and explicit
cost: $\arg\max_a [\widehat\Delta(a\mid h_1) - \lambda_B c(a)]$. Trained
and frozen on the same 100 development states
(\texttt{seq\_controller\_v2\_frozen}, sha256[:16]
\texttt{b74b8b6e70332a52}); evaluated on the already-separated 300-state
set as a disclosed secondary analysis.

\paragraph{Budget-constrained oracles.}
The Lagrangian sweep uses
$a_i(\lambda) = \arg\max_a [y_i^{(\mathrm{sel})}(a) - \lambda c_i(a)]$
over a finite $\lambda$ grid ($\lambda_{\text{tok}} \in \{0\} \cup
\mathrm{logspace}(-6,-2,33)$; $\lambda_{\text{lat}} \in \{0\} \cup
\mathrm{logspace}(-5,-0.5,33)$), selected on one replay and scored on the
other. The exact multiple-choice-knapsack DP closes the relaxation
loophole (Section~\ref{sec:results:exact}). Both are counterfactual and
non-deployable by construction.

% ======================================================================
\section{Preregistration Index and Deviations Log}
\label{sec:prereg}

All preregistration documents are included verbatim in
\texttt{preregistrations/}. Each was written and frozen before the
corresponding outcomes existed; ``post-hoc relative to the paper, but
preregistered relative to its own data'' is stated explicitly where it
applies.

\paragraph{Preregistration chronology.}
Every row lists the date the document was frozen, the SHA-256 prefix of the
frozen artifact(s) and evaluation manifest recorded \emph{in} the document,
the collection window from the run metadata of the log, the single
preregistered primary endpoint, and the outcome as reported. ``Pool'' rows
train on sets that had already been analyzed (disclosed in the document);
evaluation sets are never reused.

\begin{center}
\scriptsize
\setlength{\tabcolsep}{3pt}
\begin{tabular}{@{}>{\raggedright\arraybackslash}p{2.25cm}>{\raggedright\arraybackslash}p{0.85cm}>{\raggedright\arraybackslash}p{2.35cm}>{\raggedright\arraybackslash}p{2.15cm}>{\raggedright\arraybackslash}p{1.6cm}>{\raggedright\arraybackslash}p{2.4cm}>{\raggedright\arraybackslash}p{2.2cm}@{}}
\toprule
Document & Frozen & Frozen artifact (sha256[:16]) & Evaluation set ($n$; sha) & Collected & Primary endpoint & Outcome \\
\midrule
\texttt{FRESH\_LOCKBOX} & 08-23 & frozen confidence + v2.2 gate (unchanged) & VQAv2 250 + TextVQA 250 fresh train states & 08-23 & Format-confound replication on fresh states & Replicated (Sec.~\ref{sec:results}) \\
\texttt{MULTIACTION\_ HELDOUT} & 08-28 & multi-action policy v1 (dev-fitted) & TextVQA 250 + VQAv2 250 held-out & 08-28--30 & Multi-action $-$ best fixed at $1/4\times$, $1/2\times$ (TextVQA, strict) & PASS: $+7.0\ [+2.7,+11.4]$ / $+9.9\ [+5.5,+14.4]$ \\
\texttt{SECOND\_READER} & 08-30 & protocol only (Phi-4 reader) & Gate-0 states, 196 MME + 200 TextVQA & 08-30 & Pooled symmetric headroom $>0$ (H1) & PASS: pooled $+6.4$; TextVQA $+10.0\ [+6.0,+14.0]$ \\
\texttt{MULTI\_REPLAY} & 08-30 & protocol only (replays 3--5) & 100-state subsets, both datasets & 08-30 & Protocol validation, no threshold & $K$=2 sufficient; $K$=2 estimate conservative \\
\texttt{HUMAN\_AUDIT} & 08-28 & blinded sheets (never opened) & 381 + 489 items & --- & Strict vs.\ human label transition & \textbf{Withdrawn 08-30; declined again 09-01 and 09-20} \\
\texttt{HUMAN\_ ADJUDICATION} & 09-21 & blinded 300-item package \texttt{1384f99672bc8d65} & prospective-600 generations: 200 scorer disagreements + 100 controls & 09-21 (3 annotators) & $P(\text{human correct}\mid \text{strict}=0,\ \text{tolerant}>0)$ & Executed; protocol frozen before labels: \textbf{.807 [.746, .856]}; Fleiss' $\kappa$ .766 \\
\texttt{PROSPECTIVE\_ SELECTOR} & 08-31 & selector v2 \texttt{d4e495063ac56b59}; mixture (dev-fitted) & TextVQA 600; \texttt{a535d8edce3aa822} & 08-31 & v2 $-$ mixture, $1/2\times$ strict & \textbf{primary non-confirmatory}: $+0.9\ [-1.6,+3.4]$; lenient $+4.6\ [+2.4,+6.8]$ \\
\texttt{SEQUENTIAL} & 08-31 & protocol only (two-step audit) & TextVQA 100 & 08-31 & Sequential oracle $-$ stop-after-$a_1$ & $+9.3\ [+4.7,+14.5]$ (center) \\
\texttt{SEQ\_PROSPECTIVE} & 08-31 & seq.\ controller v1 \texttt{cffa1043f455f7d1} & TextVQA 300; \texttt{4196e1abbead1506} & 08-31 & Controller $-$ always-$\mathrm{CROP}_{\mathrm{TR}}$, center, strict & \textbf{primary non-confirmatory}: $-1.7\ [-3.6,+0.2]$ \\
\texttt{SELECTOR\_V3} & 09-01 & selector v3 (H + ridge) \texttt{4fd1fa4fc1793b7a} & TextVQA 300; \texttt{7bc5ec8a8eb9a61f} & 09-01 & v3 $-$ mixture, $1/2\times$ strict & \textbf{primary non-confirmatory}: $+1.7\ [-1.9,+5.1]$; $\lambda$-budget uncontrollable \\
\texttt{STRONG\_ CONTROLLER} & 09-20 & F0 $\times$ \{ridge, xgb, mlp\} \texttt{5c60ccab3f7d08b8} (pool 1,350) & TextVQA 400; \texttt{7312457434be55f3} & 09-20 & xgb $-$ mixture, $1/2\times$ strict & \textbf{primary non-confirmatory}: $+1.3\ [-1.4,+4.1]$; xgb $-$ ridge $-0.3\ [-2.0,+1.6]$ \\
\texttt{HIDDEN\_ CONTROLLER} & 09-20 & \{F0, H, H+F0\} $\times$ \{ridge, xgb, mlp\} \texttt{ed93f19437264e56} (pool 1,750; primary cell H/mlp) & TextVQA 400; \texttt{af41a4ac740a9ecb} & 09-21 & H/mlp $-$ mixture, $1/2\times$ strict & \textbf{primary non-confirmatory}: $+0.8\ [-2.1,+3.8]$; H/xgb $-$ F0/xgb $-3.1\ [-5.8,-0.6]$ \\
\bottomrule
\end{tabular}
\end{center}

Zero-inference analyses added in the 09-20/21 round (no preregistration
needed; existing logs only): cross-reader agreement conditional on
baseline capability, replay-pairwise robustness (all ordered replay pairs,
and the R1/R2/R3 triangle at $n$=400), and the action-matched restricted
oracle (Section~\ref{sec:results:deporacle}).

\paragraph{Original human audit withdrawal and later adjudication.}
The original blinded semantic-utility audit preregistered on 08-28 was
withdrawn on 08-30 before any human verdict was entered; its sheets remained
empty and no result from that protocol exists.

A distinct human-adjudication study was later motivated post hoc, after the
automatic scorer-sensitivity results were observed; its sampling and
annotation protocol (\texttt{PREREG\_HUMAN\_ADJUDICATION}, 09-21) were
preregistered before any human label was collected. This later study
sampled 300 unique (state, normalized answer) pairs from the 600-state
prospective TextVQA pool: 200 benchmark-negative, format-tolerant-positive
disagreement pairs and 100 agreement controls. Annotators were blinded to
scorer outputs, action arm, replay, and policy identity. This later
adjudication is treated throughout the paper as retrospective validation
rather than preregistered primary evidence.

\paragraph{Deviations and incidents (complete log).}
\begin{enumerate}
\item \textbf{Engine death during held-out collection (08-28).} A vLLM
  engine crash poisoned 183 VQAv2 held-out states with transport-error
  rollouts (the collector logged errors instead of aborting). The affected
  rollouts were stripped (backups retained) and re-collected \emph{before}
  any evaluation was run. All subsequent drivers carry a
  zero-transport-error integrity gate; every reported collection passed it.
\item \textbf{scikit-learn version shim.} Frozen sklearn-1.9.0 pickles were
  evaluated under 1.7.2 via a \texttt{predict\_proba} attribute shim,
  asserted equal to the sigmoid of the frozen coefficients to $10^{-16}$.
\item \textbf{Human audit withdrawn without execution.} The blinded human
  semantic-utility audit (preregistered 08-28) was withdrawn on 08-30
  before any verdict was entered; the blinded sheets remain empty and the
  private mapping unopened. No result exists to report or suppress. The
  paper's ``measurement sensitivity'' framing (rather than a claim about
  which scorer is semantically correct) reflects this: a separately
  proposed adjudication protocol was declined twice (09-01, 09-20) and then executed on 09-21 under a third, narrower preregistration (\texttt{HUMAN\_ADJUDICATION}; Section~\ref{sec:results:adjudication}).
\item \textbf{Manifest regeneration incident (09-01).}
  \texttt{textvqa\_prospective.jsonl} was accidentally regenerated with a
  changed exclusion pool by an old builder script; it was restored
  immediately by reconstructing the original 950-id exclusion set. The
  restored file matches the preregistered SHA-256 prefix
  \texttt{a535d8edce3aa822} exactly and its ids match the collected log
  600/600. No analysis touched the corrupted file.
\item \textbf{Phi-4 weights re-download.} The HuggingFace cache had been
  purged between sessions; weights were re-downloaded and the serving
  configuration reproduced from the original serving log.
\item \textbf{Hidden-controller pipeline relaunches (09-20/21).} The
  \texttt{HIDDEN\_CONTROLLER} pipeline was launched three times: the first
  attempt failed before any inference because the held-out set is a
  train-split draw and the extraction script was pointed at the val
  manifest (fixed); the second and third attempts lost the shared GPU to
  other users' jobs during the 1.5-hour training step and at server
  start. The frozen artifact (\texttt{ed93f19437264e56}, 09-20 15:59) was
  never retrained, no evaluation state was collected before the final
  launch (09-21 09:56), and the fresh-set hidden states were extracted
  after collection from STOP-time inputs only. The training pool
  intentionally includes every previously analyzed TextVQA set (1,750
  states); the evaluation set is disjoint from all of them.
\item \textbf{Adjudication file drop (09-21).} The first set of verdict files
  contained a duplicated file (annotator 1 and annotator 3 byte-identical).
  The preregistered single-pass analysis was not run on it; the correct
  third file was supplied and the analysis ran once on three distinct files
  (MD5 prefixes \texttt{d1ced72c}, \texttt{24f29c79}, \texttt{82d06afc}).
\item \textbf{Strong-controller collection (09-20).} Single launch; 400/400
  states, 13,200 rollouts, zero transport errors; the same log carries
  replay 3 for the replay-robustness extension.
\item \textbf{Sequential controller v2 confidence regeneration.}
  Development first-step confidence required regeneration; byte-match was
  379/400 against the original generations (disclosed in the result file).
\item \textbf{Comparator strengthening (disclosed post-hoc).} The original
  structural comparator was the best deterministic fixed arm; randomized
  state-independent mixtures were added as a disclosed post-hoc
  strengthening, and all mixture-based conclusions are labeled accordingly.
\end{enumerate}

% ======================================================================
\section{Additional Results}
\label{sec:results}

All confidence intervals resample independent perceptual states
(10{,}000 bootstrap draws); policy differences use state-paired resampling.
``pp'' denotes percentage points. Unless noted, scoring is strict.

\subsection{Budget-constrained oracle frontier (Gate-0 logs, Qwen)}

Cross-replay symmetrized, deterministic arms; budget expressed as a
fraction of the always-\textsc{HighRes} additional-token cost.

\begin{center}
\begin{tabular}{lcc}
\toprule
Budget & MME $\Delta$ vs.\ \textsc{Stop} (recovery) & TextVQA $\Delta$ (recovery) \\
\midrule
5\% & $+.122$ (96\%) & $+.112$ (59\%) \\
10\% & $+.128$ (100\%) & $+.184$ (97\%) \\
20\% & $+.128$ (100\%) & $+.190$ (100\%) \\
\bottomrule
\end{tabular}
\end{center}

\subsection{Exact constrained oracle vs.\ Lagrangian relaxation}
\label{sec:results:exact}

At $\ge$10\% budget the Lagrangian read-off matches the randomized hull to
$\le$1pp. At the 5\% budget the exact multiple-choice-knapsack DP is
\emph{higher} by $+1.5$pp (MME) and $+2.8$pp (TextVQA): closing the
relaxation loophole strengthens the front-loading conclusion. The DP's
cross-replay evaluated value can sit $\sim$1pp below the hull at loose
budgets because exact optimization of the selection replay slightly
overfits selection noise --- itself a small illustration of selection
inflation.

\subsection{Per-action damage decomposition}

$P(\text{harm} \mid \text{\textsc{Stop} correct})$, existential over the
family (``any'') vs.\ single actions, Qwen:

\begin{center}
\begin{tabular}{lcc}
\toprule
Family & MME any / single & TextVQA any / single \\
\midrule
\textsc{Crop} (grid) & .449 / .169 & .154 / .049 \\
\textsc{OCR} & .080 / .080 & .085 / .085 \\
\textsc{HighRes} & .145 / .145 & .043 / .043 \\
Any action (DMG) & .464 / .157 & .222 / .051 \\
\bottomrule
\end{tabular}
\end{center}

Aggregate DMG (13.8/13.0\%) is dominated by the grid-crop existential;
single-action harm is 4--17\%. On TextVQA,
$P(\text{help}\mid\text{\textsc{Stop} wrong}) = .782 \gg
P(\text{harm}\mid\text{\textsc{Stop} correct}) = .222$.
Phi-4 shows the same structure with lower DMG (7.7/7.0\%) but higher
\textsc{OCR} harm on MME (40\% vs.\ Qwen's 10.7\%).

\subsection{Answer-span confidence variants}
\label{sec:results:spanconf}

Regeneration byte-match 100\%. TextVQA ROC--AUC is stable at 0.74--0.80
under all span variants. VQAv2 under the original prompt: reply-mean
0.42--0.43 (below chance), recovering only to 0.58--0.62 at span-minimum;
under the one-word prompt 0.72--0.90. The confidence pathology is a
generation-format phenomenon, not a token-aggregation artifact.

\subsection{Second reader (Phi-4): lenient rescoring and realizable signal}

TextVQA lenient rescoring of the same Phi-4 generations: \textsc{Stop}
.258 $\to$ .505 (format artifact confirmed); cross-fitted best fixed .624;
cross-replay oracle .743; headroom $+11.9$pp $[+7.1, +16.7]$;
\textsc{Stop}-optimal 36.9\%; frontier recovery 61\% at 5\% budget, 90\% at
10\%. MME is multiple-choice, so lenient $=$ strict by construction.
Phi-4's \textsc{Stop} verbosity under this protocol (4.0 words mean, 32\%
$>$3 words, vs.\ 1.5 / 5\% in its own native-format runs) explains its low
absolute strict accuracy; structural quantities, not absolute levels, are
the comparanda.

Phi-4's own \textsc{Stop}-confidence selector is indistinguishable from
random allocation at every matched budget (confidence $-$ random point
estimates $-0.003$ to $+0.013$, all CIs crossing 0), while the constrained
oracle stands far above both (TextVQA 10\% budget: random .251, confidence
.254, oracle .396). Headroom replicates across readers; the zero-training
realizable signal does not.

\subsection{Multi-replay robustness (replays 3--5)}
\label{sec:results:multireplay}

On 100-state outcome-blind subsets: the naive same-replay oracle stays
inflated at every $K$ (MME .65--.66 vs.\ cross-replay .50--.53; TextVQA
.885--.888 vs.\ .773). $K$=4--5 estimates stay inside the $K$=2 bootstrap
CI on both datasets (pre-declared sufficiency criterion). TextVQA held-out
mean and symmetric estimates \emph{rise} with $K$ (.773 $\to$ .830/.836):
more evaluation replays denoise the outcome, so the published $K$=2
headroom is conservative. Headroom is positive at every $K$.

\paragraph{All ordered replay pairs, and a three-replay follow-up on 400 states.}
On the same 100-state subsets every ordered replay pair $(r \to s)$ gives a
positive cross-replay headroom (TextVQA 20/20 pairs, mean $+.110$, sd
$.035$, 16 of 20 CIs excluding zero; MME 20/20, mean $+.070$), and the
published R1$\to$R2 cell is the smallest of the 20 on the TextVQA subset. On
the later 400-state nonlinear-head follow-up, an independently constructed
third replay (``In the image shown: \ldots'') permits all six ordered
directions among replays 1--3: headroom R1$\to$R2 $+.095$ $[+.055, +.135]$,
R1$\to$R3 $+.103$ $[+.068, +.138]$, R2$\to$R3 $+.113$ $[+.082, +.147]$
(reverse directions $+.095$ / $+.102$ / $+.102$); \textsc{Stop} agreement
ranges from .945 to .980, per-arm action-score correlation from .897 to
.933, and sign agreement of the per-action utility among both-nonzero pairs
from .97 to 1.00. Files: \texttt{replay\_pairwise\_robustness},
\texttt{replay\_pairwise\_strong}.

\subsection{Randomized-mixture baselines on the held-out sets}

TextVQA strict, selector $-$ mixture (paired bootstrap):

\begin{center}
\small
\begin{tabular}{lccc}
\toprule
Budget & vs.\ in-sample & vs.\ cross-fitted & vs.\ deployable \\
\midrule
$1/4\times$ & $+.009\ [-.027,+.045]$ & $+.020\ [-.017,+.056]$ & $+.010\ [-.025,+.046]$ \\
$1/2\times$ & $+.012\ [-.022,+.046]$ & $+.018\ [-.012,+.049]$ & $+.030\ [-.003,+.063]$ \\
$1\times$   & $-.011\ [-.045,+.024]$ & $-.011\ [-.045,+.024]$ & $+.028\ [-.009,+.065]$ \\
\bottomrule
\end{tabular}
\end{center}

VQAv2 strict vs.\ cross-fitted: $-.153\ [-.204,-.102]$ at $1/4\times$ and
$-.072\ [-.117,-.027]$ at $1/2\times$ (a fixed grid crop significantly
beats the selector); under lenient rescoring the reversal vanishes
($+.010\ [-.017,+.038]$, $+.014\ [-.017,+.045]$). The deployable mixture
transferred from development to held-out with no loss (.717 development fit
$\to$ .744 held-out at $1/2\times$), so a transfer-loss explanation is not
available.

\subsection{Prospective one-step evaluation: realized-cost matching}

Re-solving the development mixture at the selector's \emph{realized} token
cost (53/129/212) on the 600 prospective states: strict
$+1.9\ [-0.1,+4.0]$ / $+1.2\ [-1.4,+3.8]$ / $+0.9\ [-1.7,+3.4]$ (all
non-confirmatory); lenient $+2.7\ [+0.9,+4.5]$ / $+5.1\ [+2.8,+7.4]$ /
$+3.6\ [+1.4,+5.9]$ (all excluding zero). Cost-matching therefore
strengthens, not weakens, the scorer-dependence finding.

\subsection{Sequential experiments: full cell table}

Development audit ($n$=100): cross-replay oracle over
\{stay, $\to\mathrm{CROP}_{\mathrm{TR}}$, $\to$\textsc{HighRes}\} beats
stopping after the first action by $+9.3$pp $[+4.7,+14.5]$ (center-crop
branch) and $+9.2$pp $[+4.2,+14.5]$ (OCR branch). Continuation gain
concentrates on first-step failures ($+22$ to $+27$pp) and is $\approx$0
on successes. ``First-step answer changed'' predicts residual utility with
AUC .70/.74 (base rate of ``second look helps'' $\approx$ 10--15\%);
answer length alone $\approx$ .55. At matched tokens the best one-step arm
(CROP-TR, .767 at 243 tokens) matches or beats every fixed two-step chain:
the sequential value is adaptive, not a fixed pipeline.

Prospective test ($n$=300), controller $-$ frozen continuation
(always-$\mathrm{CROP}_{\mathrm{TR}}$):

\begin{center}
\small
\begin{tabular}{llcc}
\toprule
Branch & Scorer & v1 (original) & v2 (cost-aware, disclosed) \\
\midrule
Center & strict & $\mathbf{-1.7\ [-3.6,+0.2]}$ (primary) & $-0.6\ [-1.8,+0.6]$ \\
Center & lenient & $-1.2\ [-2.9,+0.4]$ & --- \\
OCR & strict & $-0.3$ (ns) & $+0.4$ (ns) \\
OCR & lenient & $+2.7\ [+0.4,+5.1]$ & $+2.7\ [+0.5,+4.8]$ \\
\bottomrule
\end{tabular}
\end{center}

Both v1 variants collapse to look-rate 1.00 on fresh states; v2 restores
stopping (look-rate 0.83--0.85, or 0.30 under the conservative
$\lambda_B$, which matches always-$\mathrm{CROP}_{\mathrm{TR}}$ within
$-1.6$pp ns at 80 vs.\ 241 tokens --- one third the cost). Sequential
oracle headroom on the same 300 states (all eight CIs exclude zero):
oracle $-$ continuation $= +2.8\ [+1.1,+4.8]$ (center strict),
$+5.8\ [+3.4,+8.4]$ (OCR strict), $+2.7\ [+1.0,+4.7]$ /
$+6.2\ [+3.8,+8.8]$ lenient; oracle $-$ stop-after-first
$= +4.4\ [+1.9,+6.9]$ / $+14.0\ [+10.3,+17.8]$. Only one of eight
secondary controller cells is positive with a CI excluding zero
(multiplicity caveat in the main paper).

\subsection{Hidden-state selector v3: budget-control failure}
\label{sec:results:v3}

Primary ($1/2\times$ strict, v3 $-$ mixture): $+1.7$pp $[-1.9,+5.1]$;
$1/4\times$: $+2.9\ [-0.7,+6.3]$; $1\times$: $+1.7\ [-2.0,+5.6]$. The
pre-flagged instability materialized exactly: realized cost is
$\approx$174--179 tokens at \emph{every} nominal budget --- $\lambda(B)$
does not bind with hidden-state value heads. The $1/4\times$ and
$1/2\times$ point estimates are therefore cost-confounded (2.3$\times$ /
1.1$\times$ the mixture's tokens) and are not claimed. The only cost-fair
cell is $1\times$, where v3 is cheaper than the mixture (179 vs.\ 241
tokens) at $+1.7$pp ns strict / $+3.5\ [-0.0,+7.2]$ borderline lenient,
and v3 $\approx$ v2 ($+0.2$pp ns). Secondary v3 $-$ v2 at $1/4\times$
($+4.5\ [+1.0,+8.1]$) is likewise cost-confounded (174 vs.\ 53 tokens) and
not claimed as a win. Beyond the non-confirmatory primary, this surfaces a
deployability failure mode --- loss of $\lambda$-budget controllability ---
that the handcrafted selector does not have.

\subsection{F0 nonlinear-head prospective follow-up}
\label{sec:results:strong}

On 400 fresh TextVQA states (manifest sha256[:16] \texttt{7312457434be55f3};
13,200 rollouts including a third replay, zero transport errors) at the
strict $1/2\times$ operating point, F0/XGBoost reaches .726 at 148 visual
tokens, F0/ridge .728 at 141, and F0/MLP .716 at 160, compared with .713 for
the frozen state-independent mixture, .777 for the action-matched oracle,
and .812 for the unrestricted oracle. The preregistered XGBoost-minus-mixture
primary comparison is $+1.3$ points (95\% CI $[-1.4, +4.1]$) and therefore
non-confirmatory. Capacity contrasts are null at every budget (XGBoost
$-$ ridge $-0.3$ $[-2.0, +1.6]$, MLP $-$ ridge $-1.3$ $[-3.2, +0.6]$ at
$1/2\times$); the pool-trained ridge exceeds the published dev-only v2 by
about 3 points (.728 vs.\ .699). Realized costs track the nominal budgets for
all three heads (54--78 / 141--160 / 210--225 tokens). Secondary
format-tolerant XGBoost-minus-mixture differences are positive at
$1/2\times$ ($+3.3$ $[+0.7, +6.1]$) and $1\times$ ($+4.4$ $[+1.4, +7.5]$),
with XGBoost $-$ ridge still null. These results reduce the concern that
the original prospective result is specific to a linear ridge head, but
they do not bound richer controller classes. File:
\texttt{strong\_controller\_eval}.

\subsection{Representation--head $3\times3$ prospective test}
\label{sec:results:hidden}

On a further 400 fresh TextVQA states (manifest sha256[:16]
\texttt{af41a4ac740a9ecb}; 8,800 rollouts, zero transport errors; hidden
states extracted after collection from \textsc{Stop}-time inputs only), at
the strict $1/2\times$ operating point, accuracies (realized visual-token
cost) are:

\begin{center}
\small
\begin{tabular}{lccc}
\toprule
& F0 & H & H+F0 \\
\midrule
Ridge   & .749 (146) & .729 (163) & .732 (161) \\
XGBoost & .737 (141) & .705 (151) & .724 (131) \\
MLP     & .744 (122) & .722 (171) & .707 (142) \\
\bottomrule
\end{tabular}
\end{center}

The mixture reaches .714, the action-matched oracle .777, and the
unrestricted oracle .808. The preregistered H/MLP-minus-mixture primary
comparison is $+0.8$ points (95\% CI $[-2.1, +3.8]$) and non-confirmatory;
H/XGBoost $-$ mixture is $-0.8$ $[-3.6, +2.0]$. H/XGBoost underperforms
F0/XGBoost by $-3.1$ points ($[-5.8, -0.6]$; $-5.9$ at $1/4\times$, $-2.5$
at $1\times$) and H/ridge by $-2.4$ ($[-4.8, -0.1]$); H/ridge $-$ F0/ridge
is $-1.9$ $[-4.2, +0.4]$ and H+F0/XGBoost $-$ F0/XGBoost $-1.3$
$[-3.3, +0.7]$. Realized costs track the nominal budgets for all nine
cells (41--73 / 122--171 / 162--237 tokens), so the v3 budget-control
failure does not recur with pooled out-of-fold $\lambda$. A secondary,
unadjusted F0/ridge comparison is $+3.5$ points ($[+0.8, +6.4]$) versus
the mixture and is treated as hypothesis-generating (the same contrast was
$+1.6$ $[-1.2, +4.4]$ on the nonlinear-head set). Under format-tolerant
rescoring, all nine configurations exceed the mixture at $1/2\times$ by
$+2.8$ to $+5.2$ points with unadjusted confidence intervals excluding
zero, and the H $-$ F0 contrasts are null; these secondary analyses are not
part of the preregistered strict primary test. Pre-collection nested
cross-validation on the training pool had already ranked every H-based
cell below F0 (recorded in the frozen manifest). File:
\texttt{hidden\_controller\_eval}.

\subsection{Action-matched counterfactual oracle}
\label{sec:results:deporacle}

The unrestricted structural oracle may choose among individual crop slots
\{\textsc{Stop}, \textsc{Crop}$_{1..8}$, \textsc{OCR}, \textsc{HighRes}\},
whereas the deployed selectors choose only among
$\{\textsc{Stop},\textsc{Crop-family},\textsc{OCR},\textsc{HighRes}\}$. To
isolate the contribution of this action-space mismatch, we recompute the
cross-replay oracle over the deployed family-level action space, using the
uniform eight-crop expectation for \textsc{Crop-family} (outcome and cost
exactly as the selectors use them), under the audit protocol (select on
replay 1, ties to the cheaper family, score on replay 2) and against the
\emph{same} cross-fitted best fixed arm, so that $H^{\mathrm{full}} -
H^{\mathrm{dep}}$ isolates the action space. Both oracles remain
counterfactual because their action choice uses replay outcomes. This
analysis is post hoc (zero inference on the existing logs) and is not a
preregistered primary endpoint. Strict:

\begin{center}
\small
\begin{tabular}{lrccc}
\toprule
Set & $n$ & $H^{\mathrm{full}}$ [CI] & $H^{\mathrm{dep}}$ [CI] & matched/unrestricted \\
\midrule
TextVQA Gate-0 & 200 & $+8.9\ [+4.0,+13.9]$ & $+4.6\ [-0.1,+9.3]$ & .51 \\
MME Gate-0 & 196 & $+8.2\ [+3.3,+13.2]$ & $+3.0\ [-1.1,+7.2]$ & .37 \\
TextVQA held-out & 250 & $+7.9\ [+4.2,+11.6]$ & $+4.4\ [+1.1,+7.6]$ & .56 \\
TextVQA prospective & 600 & $+9.8\ [+7.1,+12.6]$ & $+5.6\ [+3.1,+8.2]$ & .57 \\
TextVQA v3 & 300 & $+11.4\ [+8.0,+15.1]$ & $+6.7\ [+4.0,+9.4]$ & .59 \\
TextVQA strong & 400 & $+9.4\ [+6.3,+12.5]$ & $+6.3\ [+3.6,+9.0]$ & .67 \\
TextVQA hidden & 400 & $+7.8\ [+4.7,+10.9]$ & $+4.1\ [+1.5,+6.8]$ & .53 \\
TextVQA Gate-0, Phi-4 & 200 & $+9.9\ [+5.7,+13.9]$ & $+7.5\ [+4.1,+10.8]$ & .76 \\
MME Gate-0, Phi-4 & 196 & $+2.3\ [-2.9,+7.6]$ & $+1.5\ [-3.0,+6.1]$ & .66 \\
\bottomrule
\end{tabular}
\end{center}

Restricting the oracle to the deployed family-level action space removes
roughly 40--50\% of the headroom on TextVQA (35--65\% on MME); a family-level
headroom of $+4$ to $+7$pp remains, with CIs excluding 0 on every fresh
TextVQA set ($n \ge 250$). Budget-constrained (Lagrangian) values at
$1/2\times$ are .777--.781 (action-matched) versus .809--.819
(unrestricted) on the fresh sets. At the original 600-state $1/2\times$ operating point, the
unrestricted oracle reaches .818, the action-matched oracle .777, and the
frozen selector .720: 4.1 points of the unrestricted 9.8-point
oracle--selector gap are attributable to finer crop granularity, while 5.7
points remain between the action-matched counterfactual oracle and the
realized selector. The opportunity therefore has two layers: roughly half
is crop granularity (which box), which the selectivity audit shows is not
question-predictable; the other half is family choice, which is inside the
selectors' action space and remains largely unrealized (best deployable
cell .749 vs.\ action-matched oracle .777 at $1/2\times$ on the
representation--head set). File: \texttt{deployable\_oracle}.

\subsection{Blinded human adjudication of scorer disagreement}
\label{sec:results:adjudication}

\paragraph{Sampling.}
The adjudication universe consists of the 1,584 unique (state, normalized
answer) pairs produced by the 13,200 generations of the 600-state
prospective TextVQA evaluation. We sample 200 pairs for which the benchmark
(strict) score is zero but the format-tolerant score is positive (245 such
pairs exist), with allocation proportional to the dominant action family
among the rollouts producing each pair (\textsc{Crop} 88, \textsc{HighRes}
49, \textsc{OCR} 42, \textsc{Stop} 21). We additionally sample 100 agreement
controls: 50 pairs scored correct by both automatic scorers and 50 scored
incorrect by both. Sampling is fixed by preregistration (seed 20260921)
before any human label exists; the 300 pairs cover 1,527 rollouts, the 200
disagreement pairs 675.

\paragraph{Blinding and adjudication.}
Three independent annotators judge semantic correctness (correct /
incorrect / ambiguous) from the image, the question, the reference answers
with vote counts, and the model prediction only, without access to
automatic scorer outputs, action arm, replay identity, or policy identity,
under the instruction that containing the reference token is not sufficient
for correctness. Majority vote defines the adjudicated label; a majority of
``ambiguous'' and pairs without a majority are retained separately.

\paragraph{Agreement and semantic error.}
Across the 300 sampled pairs, raw pairwise agreement is .906 (mean; Cohen's
$\kappa$ .70 / .88 / .72), Fleiss' $\kappa=.766$, and Krippendorff's
$\alpha=.767$; 292 pairs receive a majority label (258 unanimous), 7 are
ambiguous by majority, 1 has no majority. Among the 197 resolved
benchmark-negative, format-tolerant-positive pairs, 80.7\% (Wilson 95\% CI
$[74.6, 85.6]$) are human-correct and 19.3\% ($[14.4, 25.4]$) are
human-incorrect. The human-incorrect share concentrates in partial-credit
matches: 12.3\% of pairs with format-tolerant score 1 ($n$=154) versus
43--46\% at scores $1/3$ and $2/3$ ($n$=43). All 50 agreement-control pairs
scored correct by both automatic scorers are human-correct; among the 45
resolved pairs scored incorrect by both, 11 are judged human-correct
(paraphrases neither scorer accepts). Weighting the three strata by their
population sizes (245 / 580 / 670), agreement with the human label is
.758 $[.699, .812]$ for the benchmark scorer and .856 $[.797, .909]$ for the
format-tolerant rule.

\paragraph{Policy-level sensitivity analysis.}
For every rollout whose (state, answer) pair was adjudicated with a
majority label, the human label replaces the automatic score without
reweighting (659 rollouts); all other rollouts retain their automatic
scores. For unresolved or unsampled scorer-disagreement pairs, the
pessimistic assumption retains the benchmark score of zero, whereas the
optimistic assumption treats only format-tolerant full-credit matches as
correct. All partial-credit matches, including scores of $1/3$ and $2/3$,
remain incorrect under both assumptions. These are sensitivity bounds over unresolved
scorer-disagreement cases, not formal bounds on semantic accuracy, because
unadjudicated agreement cases remain automatically scored.

At the primary $1/2\times$ operating point, frozen selector minus frozen
mixture changes from $+0.9$ points ($[-1.6, +3.4]$) under benchmark scoring
to $+4.1$ ($[+1.8, +6.5]$) and $+4.2$ ($[+2.0, +6.5]$) points under the
pessimistic and optimistic assumptions; the format-tolerant estimate is
$+4.6$ ($[+2.4, +6.8]$). At $1/4\times$ / $1\times$ the pessimistic values
are $+2.5$ $[+0.2, +4.9]$ / $+3.0$ $[+0.5, +5.5]$ and the optimistic values
$+2.7$ $[+0.4, +5.0]$ / $+2.9$ $[+0.5, +5.4]$. Only policies frozen
before the 600-state set was collected are evaluated here. File:
\texttt{human\_adjudication}.

\subsection{Cross-dataset and cross-reader transfer details}

\paragraph{Cross-dataset (TextVQA selector $\to$ VQAv2, zero adaptation).}
Strict: parity with the cross-fitted mixture ($-0.9$/$-2.7$/$-1.9$pp, CIs
include 0); the selector loses to the in-sample mixture at
$1/2\times$/$1\times$ (that bound is baseline-favoring by construction).
Lenient: directionally positive at every budget ($+2.8$/$+2.0$/$+1.5$pp,
CIs include 0). The selector's action mix shifts sensibly (OCR-heavy on
TextVQA $\to$ \textsc{Crop}/\textsc{HighRes}-heavy on VQAv2). Note: the
oracle column in this analysis is the family-level allocation ceiling
(4 arms, crop = uniform expectation) while mixtures draw on per-slot arms.

\paragraph{Cross-reader (Qwen selector $\to$ Phi-4 states).}
MME: the transferred selector nearly always stops (realized cost
$\le$73 tokens) and loses to the target-reader cross-fitted mixture by
$-8.0$ to $-10.4$pp (CIs exclude 0), leaving Phi-4's headroom unrealized.
TextVQA: the transferred selector beats always-\textsc{Stop} by $+7.9$ to
$+10.4$pp (all CIs exclude 0) and is directionally ahead of the
cross-fitted mixture ($+1.3$ to $+2.7$pp, CIs marginal). A
confidence-recalibrated variant barely differs: the issue is the signal,
not its calibration. Combined with the instance-level agreement analysis
(help-set overlap $\approx$ chance on TextVQA; best-family agreement among
both-help states 73--77\%), what transfers is a coarse task-level signal of
which action class helps, not instance-level perception need.

\paragraph{Agreement conditional on baseline capability.}
Because the readers' \textsc{Stop} accuracies differ widely, the pooled
help-set overlap partly reflects capability rather than perception need.
Stratifying the Gate-0 TextVQA states by whether each reader's \textsc{Stop}
is correct: among both-\textsc{Stop}-wrong states ($n$=70 strict / 51
lenient) the Qwen--Phi-4 help Jaccard rises from the pooled .19 to .35
$[.23, .47]$ (lenient .64 $[.50, .77]$) and the correlation of per-state
best-arm utility to .31 $[.14, .47]$ (lenient .48 $[.25, .69]$), while
chance-corrected $\kappa$ stays $\approx .1$ because both readers have high
help rates there; the mixed stratum is anti-correlated ($-.29$
$[-.40, -.17]$). MME both-wrong: Jaccard .40, $\kappa$ .36 $[.17, .54]$.
The instance-level claim is therefore ``moderate overlap, weak
beyond-base-rate agreement'', not independence. File:
\texttt{cross\_reader\_conditional}.

\subsection{Latency at the token-matched operating points}

Already-logged latencies of the same frozen operating points (never
recomputed, never summed with tokens): the multi-action policy's low-token
operating points substitute OCR and carry its tool time --- TextVQA
$1/4\times$: 29.0\,s/state vs.\ 0.15\,s for the confidence binary at the
same token cost; VQAv2 $1/4\times$: 11.1\,s vs.\ 0.14\,s; always-\textsc{HighRes}
0.67\,s. On the latency axis the $\lambda_{\text{lat}}$ sweep avoids OCR
and the binaries reach always-\textsc{HighRes} cheaply; OCR-heavy points
sit 1--2 orders of magnitude to the right. Any efficiency claim must name
its axis.

% ======================================================================
\section{Infrastructure and Reproducibility Details}
\label{sec:infra}

\paragraph{Models and serving.}
Qwen2.5-VL-7B-Instruct and Phi-4-multimodal-instruct, served with vLLM
0.11.0 (torch 2.8.0, CUDA 12.8), greedy decoding, fixed seed, logprobs
enabled for \textsc{Stop} rollouts. Prefix caching is disabled ---
mandatory, because the protocol issues many near-identical requests per
image and cached prefixes would couple arms that must be independent.
Model weights are frozen throughout; no fine-tuning anywhere.

\paragraph{Determinism and caching.}
All generations are cached in SQLite keyed by model, prompt version,
decoding parameters, and view bytes; changing any of these invalidates the
cache rather than silently mixing configurations. \textsc{Stop}
re-generation byte-match checks are asserted at collection time (0
mismatches on all reported collections; one disclosed exception: the
sequential-v2 development confidence regeneration at 379/400).

\paragraph{Scale.}
In addition to the structural audits (7,000 rollouts per reader and
dataset) and the held-out, lockbox, and answer-format collections, the
prospective program includes the 600-state one-step selector evaluation
(13,200 rollouts), the 300-state sequential evaluation (5,400), the
300-state hidden-state evaluation (6,600), the 400-state nonlinear-head
evaluation (13,200, including a third replay), and the 400-state
representation--head evaluation (8,800), plus 1,650 hidden-state
extractions. In total the released logs hold about 117,000 model rollouts
(excluding offline rescoring files and smoke runs). The human adjudication
covers 300 unique state--answer pairs sampled from the 600-state prospective
TextVQA collection. Every model-based evaluation collection passed the
zero-transport-error integrity gate.

\paragraph{Sampling and leakage control.}
Every evaluation set is drawn with a fixed seed from a pool that excludes
all previously used ids (a cumulative ``burn list''); manifest SHA-256
prefixes are recorded in the corresponding preregistration document before
collection and verified in the released manifests
(the 600-state prospective manifest \texttt{a535d8edce3aa822}; the
sequential 300-state manifest \texttt{4196e1abbead1506}; the v3 manifest
\texttt{7bc5ec8a8eb9a61f}; the nonlinear-head 400-state manifest
\texttt{7312457434be55f3}; the representation--head 400-state manifest
\texttt{af41a4ac740a9ecb}; the adjudication sample mapping
\texttt{1384f99672bc8d65}).

\paragraph{Prospective-set separation and sequential development.}
The one-step selector (600), sequential (300), hidden-state (300),
nonlinear-head (400), and representation--head (400) prospective evaluation
sets are mutually disjoint, disjoint from the Gate-0 development set (200,
val split) and from the train-split development, lockbox, and held-out sets
(250 each); pairwise overlap is zero for all nine sets. Each controller is
evaluated only on a set that was never used for its own training or model
selection. The research program is sequential rather than a collection of
statistically independent replications: later controllers use previously
analyzed sets in their training pools (1,350 and 1,750 states), but are
evaluated on newly drawn untouched sets. Accordingly, the paper treats the
preregistered primary endpoints as distinct prospective follow-ups and makes
no familywise-confirmatory omnibus claim across them.

\paragraph{Statistics.}
State-level bootstrap (10{,}000 draws) over independent perceptual states;
paired resampling for policy differences; no rollout-level resampling
anywhere. Preregistered primary endpoints are single comparisons; all
secondary comparisons are reported unadjusted and labeled as such.

% ======================================================================
\section{Contents of the Supplementary ZIP}
\label{sec:zip}

\begin{description}[style=nextline]
\item[\texttt{supplementary.pdf}] This document.
\item[\texttt{code/}] The complete audit framework and experiment scripts:
  \texttt{ppu/} (the library: budget accounting, actions, regions, OCR,
  prompts, replays, evaluator, cache, runner, audit), \texttt{scripts/}
  (every collection, analysis, and figure script named in the result
  files), \texttt{tests/} (offline tests exercising the full pipeline with
  stub readers --- no GPU or data required), and a README with environment
  and reproduction instructions. Absolute paths have been replaced by
  placeholders for anonymity.
\item[\texttt{preregistrations/}] All preregistration documents used in
  the study, verbatim, including the withdrawn original human-audit
  preregistration and the later preregistered prospective follow-ups
  (nonlinear-head, representation--head) and human-adjudication protocol.
\item[\texttt{results/}] Machine-readable result files (JSON) and their
  human-readable summaries (Markdown) for every analysis cited in the
  paper or in this document --- including the action-matched oracle
  analyses, the human-adjudication summaries (blinded sample mapping, the
  three verdict files, and the analysis output), the nonlinear-head
  follow-up, the representation--head grid, the replay-pairwise and
  conditional cross-reader analyses --- plus per-state CSV rows for the
  sequential prospective oracle and history analyses, and the frozen
  results digest.
\item[\texttt{models/}] Freeze artifacts for every learned component
  (pickle + JSON manifest with weights, scaler, confidence ECDF, feature
  list, and seed), matching the SHA-256 prefixes cited above.
\item[\texttt{logs/}] Complete raw rollout logs (gzipped JSONL) for every
  core experiment: the Gate-0 audits (both readers), the 600-state one-step
  prospective evaluation, the sequential development and 300-state
  prospective collections, the v3 prospective collection, the 400-state
  nonlinear-head and representation--head prospective collections, the
  held-out multi-action arms, the multi-replay robustness collections, the
  answer-format intervention arms, and the sealed binary-utility sets.
  Each row stores the full prompt metadata, served question, view labels
  and realized token costs, the \emph{raw generation}, extracted answer,
  per-metric score, and (for \textsc{Stop}) token log-probabilities ---
  every number in the paper is recomputable from these files without GPU
  inference, including under alternative scoring rules.
\item[\texttt{manifests/}] Evaluation-set manifests (sample ids,
  questions, references, native sizes, relative image paths) for all
  prospective, held-out, lockbox, and fresh sets (including the
  nonlinear-head and representation--head 400-state sets and the frozen
  300-pair adjudication sample), verifying the preregistered SHA-256
  prefixes. Images themselves are from the public
  TextVQA, VQAv2, and MME-RealWorld-Lite releases and are not
  redistributed.
\end{description}

\paragraph{Reproduction.}
Scoring, audit, and policy analyses are reproducible offline from
\texttt{results/} and \texttt{manifests/} without GPU inference (raw
generations are embedded in the result JSONs where needed; the evaluator
is deterministic). Full re-collection requires serving the two readers
with vLLM (scripts provided; prefix caching must remain disabled) and the
public benchmark images. The offline test suite
(\texttt{python -m pytest tests -q}) exercises every code path with stub
readers and stub OCR.

\end{document}
